\documentclass[10pt,letterpaper,sans]{moderncv}
\moderncvstyle{banking}
\moderncvcolor{blue}

\usepackage[scale=0.84]{geometry}
\nopagenumbers{}
\setlength{\hintscolumnwidth}{2.35cm}

\name{Krishnaa}{Vadivel}
\title{Research Engineer | Scientific Computing | Machine Learning | Electrical Engineering}
\email{srikrishnaa.careers@gmail.com}
\phone[mobile]{516-853-5663}
\social[linkedin]{sri-krishnaa-ece-phd}
\extrainfo{\href{https://krishnaa423.github.io/personal-website/}{Website} \textbar{} \href{https://github.com/krishnaa423}{GitHub: krishnaa423} \textbar{} \href{https://leetcode.com/u/krishnaa42342}{LeetCode: krishnaa42342} \textbar{} \href{https://hub.docker.com/u/krishnaa42342}{Docker Hub: krishnaa42342}}

\begin{document}
\makecvtitle

\section{Profile}
\cvitem{}{Broad-profile electrical engineering PhD researcher with training and project work spanning computational physics, semiconductor physics, integrated photonics, scientific machine learning, embedded systems, digital design, and high-performance scientific software. The core strength is connecting mathematical models, physical intuition, and implementation detail across research and engineering settings.}

\section{Research Experience}
\cventry{2021--present}{Research Assistant / PhD Researcher}{Yale University, Electrical Engineering}{}{}{\begin{itemize}%
\item Developed graph neural networks and automatic-differentiation models in Python to predict material properties including optical emission spectra.
\item Developed first-principles and many-body simulation software in Python and C++ using MPI, OpenMP, PETSc, SLEPc, CUDA, ROCm, and CMake for large simulation and linear-algebra workloads.
\item Built distributed and accelerator-enabled applications for Frontier, Perlmutter, Frontera, and related clusters, spanning both NVIDIA and AMD GPU systems.
\item Worked on light-matter interaction and quantum many-body modeling, including self-trapped-exciton and exciton-polaron problems relevant to semiconductor and photonic materials.
\item Developed early-PhD cleanroom and characterization experience through GaN power-device fabrication, MOCVD, selective-area etching, SEM, cathodoluminescence, and Raman analysis.
\item Built agentic automation with LangChain, OpenAI API models, and local Ollama models to read documentation and papers, then generate structured inputs for first-principles software.
\item Support embedded-systems instruction and hardware-oriented lab work through teaching-assistant responsibilities tied to firmware, debugging, and instrumentation.
\item Co-authored a 2025 \textit{Journal of the American Chemical Society} paper on ultrafast self-trapped exciton formation in lead-free double halide perovskites and presented related APS talks in 2024, 2025, and 2026.
\end{itemize}}

\section{Professional Experience}
\cventry{2019--2021}{Software and Embedded Systems Engineer}{Probe Technologies Inc. / Weatherford Wireline Services}{}{}{\begin{itemize}%
\item Built parallel applications with CUDA and OpenCL fallback for large acoustic waveform and sensor datasets used in engineering diagnostics and field-tool analysis.
\item Developed embedded software and digital logic for sensing, logging, and control functions using Lattice FPGA, PIC32, dsPIC, STM32, and C8051-based systems.
\item Supported board- and system-level engineering tasks including oscilloscope-based debugging, integration testing, field validation, and selective PCB modifications in Altium.
\item Developed full-stack server applications with ASP.NET and Microsoft SQL Server for real-time control of company tool testing, data collection, and data processing.
\item Developed C\# GUI applications for preparing customer inputs and interacting with deployed hardware systems.
\item Used \texttt{dolfinx}-based finite-element analysis for frequency analysis of steel pipes in wells.
\end{itemize}}
\cventry{2018--2019}{Photonics Technical Writer}{FindLight}{}{}{\begin{itemize}%
\item Produced technical articles and engineering-facing content covering lasers, optics, photonics components, and optoelectronic systems.
\end{itemize}}

\section{Selected Projects}
\cventry{}{Machine Learning Models for Material Properties}{}{}{}{\begin{itemize}%
\item Developed PyTorch Geometric graph neural networks and automatic-differentiation models for molecular-property prediction and optical-properties calculations in material systems.
\end{itemize}}
\cventry{}{Semiconductor Device Simulation}{}{}{}{\begin{itemize}%
\item Built numerical treatments of electrostatics and transport equations, including Poisson, drift-diffusion, and pn-junction-style models linking semiconductor theory with simulation practice.
\end{itemize}}
\cventry{}{Integrated Photonics Bandpass Filter}{}{}{}{\begin{itemize}%
\item Developed Meep FDTD simulations for a ring-resonator-based integrated-photonics bandpass filter that was later fabricated and tested with an erbium-doped fiber laser.
\end{itemize}}
\cventry{}{Container Orchestration for First-Principles Computing}{}{}{}{\begin{itemize}%
\item Developed CUDA-ready and cross-environment containers for first-principles software stacks and published them to Docker Hub for portable scientific computing.
\end{itemize}}
\cventry{}{FPGA AI Accelerator and Digital Design Work}{}{}{}{\begin{itemize}%
\item Implemented Verilog-based accelerator and processor projects to demonstrate datapath design, verification workflows, and hardware-oriented compute thinking.
\end{itemize}}
\cventry{}{Embedded Robotics and Human-Interface Projects}{}{}{}{\begin{itemize}%
\item Built robotics and embedded-system class projects including a surface-EMG robotic hand that was later published in Circuit Cellar, plus maze-navigation work with wireless coordination.
\end{itemize}}
\cventry{}{Integrated Photonics and FDTD Modeling}{}{}{}{\begin{itemize}%
\item Simulated band-pass and resonator-style photonic structures with Meep-based FDTD workflows for integrated-optics exploration.
\end{itemize}}
\cventry{}{Agentic Scientific Automation}{}{}{}{\begin{itemize}%
\item Built retrieval-driven tooling that reads software documentation and research papers to generate structured inputs and workflow scaffolds for computational science.
\end{itemize}}

\section{Education}
\cventry{2021--present}{PhD, Electrical Engineering}{Yale University}{}{}{Ongoing work spanning computational materials, many-body theory, semiconductor-relevant modeling, and scientific software.}
\cventry{2023}{MPhil, Electrical Engineering}{Yale University}{}{}{}
\cventry{2023}{MS, Electrical Engineering}{Yale University}{}{}{}
\cventry{2017--2018}{MEng, Electrical and Computer Engineering}{Cornell University}{}{}{Included stochastic-computing-related work, device-oriented EE topics, and advanced design coursework.}
\cventry{2013--2017}{BS, Electrical and Computer Engineering}{Cornell University}{}{}{Included embedded systems, robotics, controls, and strong math/physics/computing preparation.}

\section{Selected Publications}
\cvitem{2025}{de Paula, A. M., Li, S., Hou, B., Vadivel, S., Teles-Ferreira, D. C., Iudica, A., Kabacinski, P., Hosseini, H., McArthur, J., Cerullo, G., et al. (2025). Time-domain observation of ultrafast self-trapped exciton formation in lead-free double halide perovskites. \textit{Journal of the American Chemical Society, 147}(32), 28923--28931.}

\section{Selected Conference Presentations}
\cvitem{2026}{Vadivel, S., Grande, R. D., Strubbe, D. A., \& Qiu, D. Y. (2026). First-principles study of exciton-polarons and self-trapped excitons. \textit{APS Global Physics Summit 2026}.}
\cvitem{2025}{Vadivel, S., Qiu, D. Y., Grande, R. D., \& Strubbe, D. A. (2025). First-principles study of self-trapped exciton formation in double perovskites. \textit{APS Global Physics Summit 2025}.}
\cvitem{2024}{Vadivel, S., \& Qiu, D. (2024). First-principles study of self-trapped exciton formation in double perovskites using excited state forces. \textit{APS March Meeting Abstracts 2024}, F01--007.}

\section{Selected Coursework}
\cvitem{Mathematics}{Differential Geometry; Abstract Algebra; Honors Mathematical Analysis II}
\cvitem{Computer Science}{Theoretical Computer Science; Probabilistic Machine Learning}
\cvitem{Physics}{Graduate Quantum Mechanics I--II; Solid State Physics I--II; Many-Body Quantum Physics}
\cvitem{Electrical Engineering}{Semiconductor Physics; Thin Film Science; Lasers and Optoelectronics; Integrated Photonics; VLSI Design; Analog VLSI Design; Mathematics of Signals and Systems; Advanced Embedded Systems}

\section{Technical Skills}
\cvitem{Languages}{Python, C, C++, CUDA, JavaScript, Bash, PowerShell, C\#}
\cvitem{Scientific Computing}{MPI, OpenMP, PETSc, SLEPc, mpi4py, petsc4py, slepc4py, NumPy, pandas, SciPy, SymPy, matplotlib}
\cvitem{Machine Learning}{PyTorch, PyTorch Geometric, scikit-learn, graph neural networks, transformers, agentic automation}
\cvitem{Hardware}{Verilog, Lattice FPGA, PIC32, dsPIC, STM32, C8051, sensor interfaces, PCB bring-up, oscilloscope debugging}
\cvitem{Software / Systems}{ASP.NET, Microsoft SQL Server, Linux command line, Docker, Docker Compose, Kubernetes, gcloud}
\cvitem{HPC Platforms}{Frontier, Frontera, Perlmutter, and related national-lab or university HPC environments}

\end{document}
